
\documentclass{article}
\usepackage{colortbl}
\usepackage{makecell}
\usepackage{multirow}
\usepackage{supertabular}

\begin{document}

\newcounter{utterance}

\centering \large Interaction Transcript for game `matchit\_ascii', experiment `same\_grid', episode 1 with gemma{-}4{-}26b{-}a4b{-}it.
\vspace{24pt}

{ \footnotesize  \setcounter{utterance}{1}
\setlength{\tabcolsep}{0pt}
\begin{supertabular}{c@{$\;$}|p{.15\linewidth}@{}p{.15\linewidth}p{.15\linewidth}p{.15\linewidth}p{.15\linewidth}p{.15\linewidth}}
    \# & $\;$A & \multicolumn{4}{c}{Game Master} & $\;\:$B\\
    \hline

    \theutterance \stepcounter{utterance}  
    & & \multicolumn{4}{p{0.6\linewidth}}{
        \cellcolor[rgb]{0.9,0.9,0.9}{
            \makecell[{{p{\linewidth}}}]{
                \texttt{\tiny{[P1$\langle$GM]}}
                \texttt{Te osalete ühises arvamismängus.} \\
\texttt{Eesmärk on välja selgitada, kas see ruudustik ja teine pilt, mida ainult mina näen, on samad.} \\
\texttt{Palun kirjeldage kõigepealt oma ruudustikku. Seejärel annan mina oma kirjelduse ja me saame teineteisele esitada küsimusi ruudustike kohta, et välja selgitada, kas need on samad. Alustage nüüd oma lühikest ruudustiku kirjeldust märgiga "KIRJELDUS:", millele järgneb kirjeldus. Ärge lisage midagi muud. See on teie ruudustik:} \\
\texttt{▢ ▢ X ▢ ▢} \\
\texttt{▢ ▢ X ▢ ▢} \\
\texttt{▢ ▢ X ▢ ▢} \\
\texttt{▢ ▢ X ▢ ▢} \\
\texttt{X X X X X} \\
            }
        }
    }
    & \\ \\

    \theutterance \stepcounter{utterance}  
    & \multicolumn{4}{p{0.6\linewidth}}{
        \cellcolor[rgb]{0.9,0.9,0.9}{
            \makecell[{{p{\linewidth}}}]{
                \texttt{\tiny{[P1$\rangle$GM]}}
                \texttt{KIRJELDUS: 5x5 ruudustik, kus on järgmine muster: kolm vertikaalset X{-}i kolmandas veerus (ruudud r2c3, r3c3, r4c3, r5c3) ja üks horisontaalne rida X{-}e viimasest reast (r5c1, r5c2, r5c3, r5c4, r5c5). Täpsemalt: kolmas veerg on X{-}i fool, mis algab teisest reast ja lõpeb viimasest reast, ning viimane rida on täis X{-}e.} \\
            }
        }
    }
    & & \\ \\

    \theutterance \stepcounter{utterance}  
    & & & \multicolumn{2}{p{0.3\linewidth}}{
        \cellcolor[rgb]{0.9,0.9,0.9}{
            \makecell[{{p{\linewidth}}}]{
                \texttt{\tiny{[GM$|$GM]}}
                \texttt{continue} \\
            }
        }
    }
    & & \\ \\

    \theutterance \stepcounter{utterance}  
    & & \multicolumn{4}{p{0.6\linewidth}}{
        \cellcolor[rgb]{0.9,0.9,0.9}{
            \makecell[{{p{\linewidth}}}]{
                \texttt{\tiny{[GM$\rangle$P2]}}
                \texttt{Te osalete ühises arvamismängus.} \\
\texttt{Eesmärk on välja selgitada, kas see ruudustik ja teine pilt, mida ainult mina näen, on samad.} \\
\texttt{Palun kirjeldage kõigepealt oma ruudustikku. Seejärel annan mina oma kirjelduse ja me saame teineteisele esitada küsimusi ruudustike kohta, et välja selgitada, kas need on samad. Alustage nüüd oma lühikest ruudustiku kirjeldust märgiga "KIRJELDUS:", millele järgneb kirjeldus. Ärge lisage midagi muud. See on teie ruudustik:} \\
\texttt{▢ ▢ X ▢ ▢} \\
\texttt{▢ ▢ X ▢ ▢} \\
\texttt{▢ ▢ X ▢ ▢} \\
\texttt{▢ ▢ X ▢ ▢} \\
\texttt{X X X X X} \\
            }
        }
    }
    & \\ \\

    \theutterance \stepcounter{utterance}  
    & & & \multicolumn{4}{p{0.6\linewidth}}{
        \cellcolor[rgb]{0.9,0.9,0.9}{
            \makecell[{{p{\linewidth}}}]{
                \texttt{\tiny{[GM$\langle$P2]}}
                \texttt{KIRJELDUS: 5x5 ruudustik, kus on vertikaalne rida X{-}i kolmas veerg (ruudud 3, 8, 13, 18) ja horisontaalne rida X{-}i viies rida (ruudud 21, 22, 23, 24, 25).} \\
            }
        }
    }
     \\ \\

    \theutterance \stepcounter{utterance}  
    & & & \multicolumn{2}{p{0.3\linewidth}}{
        \cellcolor[rgb]{0.9,0.9,0.9}{
            \makecell[{{p{\linewidth}}}]{
                \texttt{\tiny{[GM$|$GM]}}
                \texttt{continue} \\
            }
        }
    }
    & & \\ \\

    \theutterance \stepcounter{utterance}  
    & & \multicolumn{4}{p{0.6\linewidth}}{
        \cellcolor[rgb]{0.9,0.9,0.9}{
            \makecell[{{p{\linewidth}}}]{
                \texttt{\tiny{[GM$\rangle$P2]}}
                \texttt{See on minu KIRJELDUSKIRJELDUS: 5x5 ruudustik, kus on järgmine muster: kolm vertikaalset X{-}i kolmandas veerus (ruudud r2c3, r3c3, r4c3, r5c3) ja üks horisontaalne rida X{-}e viimasest reast (r5c1, r5c2, r5c3, r5c4, r5c5). Täpsemalt: kolmas veerg on X{-}i fool, mis algab teisest reast ja lõpeb viimasest reast, ning viimane rida on täis X{-}e.} \\
\texttt{Nüüd esitage küsimus, et teada saada minu ruudustiku uusi aspekte, mis võivad erineda teie ruudustikust. Alustage "KÜSIMUS:" ja ärge lisage midagi muud.} \\
            }
        }
    }
    & \\ \\

    \theutterance \stepcounter{utterance}  
    & & & \multicolumn{4}{p{0.6\linewidth}}{
        \cellcolor[rgb]{0.9,0.9,0.9}{
            \makecell[{{p{\linewidth}}}]{
                \texttt{\tiny{[GM$\langle$P2]}}
                \texttt{KÜSIMUS: Kas teie ruudustikus on muid X{-}i sümboleid peale teie nimetatud vertikaalse ja horisontaalne rida, või on need kõik ainus olemasolev muster?} \\
            }
        }
    }
     \\ \\

    \theutterance \stepcounter{utterance}  
    & & & \multicolumn{2}{p{0.3\linewidth}}{
        \cellcolor[rgb]{0.9,0.9,0.9}{
            \makecell[{{p{\linewidth}}}]{
                \texttt{\tiny{[GM$|$GM]}}
                \texttt{continue} \\
            }
        }
    }
    & & \\ \\

    \theutterance \stepcounter{utterance}  
    & & \multicolumn{4}{p{0.6\linewidth}}{
        \cellcolor[rgb]{0.9,0.9,0.9}{
            \makecell[{{p{\linewidth}}}]{
                \texttt{\tiny{[P1$\langle$GM]}}
                \texttt{See on minu KIRJELDUSKIRJELDUS: 5x5 ruudustik, kus on vertikaalne rida X{-}i kolmas veerg (ruudud 3, 8, 13, 18) ja horisontaalne rida X{-}i viies rida (ruudud 21, 22, 23, 24, 25).} \\
\texttt{KÜSIMUS: Kas teie ruudustikus on muid X{-}i sümboleid peale teie nimetatud vertikaalse ja horisontaalne rida, või on need kõik ainus olemasolev muster?Alustage oma vastust "VASTUS:" ja ärge lisage midagi muud.} \\
            }
        }
    }
    & \\ \\

    \theutterance \stepcounter{utterance}  
    & \multicolumn{4}{p{0.6\linewidth}}{
        \cellcolor[rgb]{0.9,0.9,0.9}{
            \makecell[{{p{\linewidth}}}]{
                \texttt{\tiny{[P1$\rangle$GM]}}
                \texttt{VASTUS: Need on kõik ainus olemasolev muster. Minu kirjelduses on need samad elemendid, mis teie kirjelduses, kuigi me kasutasime erinevaid ruudustiku tähistusi (minu r2c3 on teie ruut 8, minu r3c3 on teie ruut 13 jne).} \\
            }
        }
    }
    & & \\ \\

    \theutterance \stepcounter{utterance}  
    & & & \multicolumn{2}{p{0.3\linewidth}}{
        \cellcolor[rgb]{0.9,0.9,0.9}{
            \makecell[{{p{\linewidth}}}]{
                \texttt{\tiny{[GM$|$GM]}}
                \texttt{continue} \\
            }
        }
    }
    & & \\ \\

    \theutterance \stepcounter{utterance}  
    & & \multicolumn{4}{p{0.6\linewidth}}{
        \cellcolor[rgb]{0.9,0.9,0.9}{
            \makecell[{{p{\linewidth}}}]{
                \texttt{\tiny{[P1$\langle$GM]}}
                \texttt{Nüüd esitage küsimus, et teada saada minu ruudustiku uusi aspekte, mis võivad erineda teie ruudustikust. Alustage "KÜSIMUS:" ja ärge lisage midagi muud.} \\
            }
        }
    }
    & \\ \\

    \theutterance \stepcounter{utterance}  
    & \multicolumn{4}{p{0.6\linewidth}}{
        \cellcolor[rgb]{0.9,0.9,0.9}{
            \makecell[{{p{\linewidth}}}]{
                \texttt{\tiny{[P1$\rangle$GM]}}
                \texttt{KÜSIMUS: Kas teie vertikaalne rida algab esimesest reast (ruut 3) või teisest reast (ruut 8)?} \\
            }
        }
    }
    & & \\ \\

    \theutterance \stepcounter{utterance}  
    & & & \multicolumn{2}{p{0.3\linewidth}}{
        \cellcolor[rgb]{0.9,0.9,0.9}{
            \makecell[{{p{\linewidth}}}]{
                \texttt{\tiny{[GM$|$GM]}}
                \texttt{continue} \\
            }
        }
    }
    & & \\ \\

    \theutterance \stepcounter{utterance}  
    & & \multicolumn{4}{p{0.6\linewidth}}{
        \cellcolor[rgb]{0.9,0.9,0.9}{
            \makecell[{{p{\linewidth}}}]{
                \texttt{\tiny{[GM$\rangle$P2]}}
                \texttt{VASTUS: Need on kõik ainus olemasolev muster. Minu kirjelduses on need samad elemendid, mis teie kirjelduses, kuigi me kasutasime erinevaid ruudustiku tähistusi (minu r2c3 on teie ruut 8, minu r3c3 on teie ruut 13 jne).} \\
\texttt{KÜSIMUS: Kas teie vertikaalne rida algab esimesest reast (ruut 3) või teisest reast (ruut 8)?Alustage oma vastust "VASTUS:" ja ärge lisage midagi muud.} \\
            }
        }
    }
    & \\ \\

    \theutterance \stepcounter{utterance}  
    & & & \multicolumn{4}{p{0.6\linewidth}}{
        \cellcolor[rgb]{0.9,0.9,0.9}{
            \makecell[{{p{\linewidth}}}]{
                \texttt{\tiny{[GM$\langle$P2]}}
                \texttt{VASTUS: Minu vertikaalne rida algab teisest reast (ruut 8). Seega meie ruudustikud ei ole täpselt samad, kuna minu vertikaalne rida on teie vertikaalse rida, kuid minu esimene rida (ruut 3) on tühi, samas kui teie kirjeldus viitab sellele, et teie vertikaalne rida algab teisest reast (mis ongi sama). Oota, kontrollin uuesti: minu kirjeldus ütles "ruudud 3, 8, 13, 18", mis tähendab esimest rida, teist, kolmandat ja neljandat. Teie kirjeldus aga ütles, et see algab teisest reast. Seega minu ruudustik on teie ruudustikust erinev, sest minu vertikaalne rida on üks ruut pikem (algab ruudust 3).} \\
            }
        }
    }
     \\ \\

    \theutterance \stepcounter{utterance}  
    & & & \multicolumn{2}{p{0.3\linewidth}}{
        \cellcolor[rgb]{0.9,0.9,0.9}{
            \makecell[{{p{\linewidth}}}]{
                \texttt{\tiny{[GM$|$GM]}}
                \texttt{continue} \\
            }
        }
    }
    & & \\ \\

    \theutterance \stepcounter{utterance}  
    & & \multicolumn{4}{p{0.6\linewidth}}{
        \cellcolor[rgb]{0.9,0.9,0.9}{
            \makecell[{{p{\linewidth}}}]{
                \texttt{\tiny{[GM$\rangle$P2]}}
                \texttt{Nüüd esitage küsimus, et teada saada minu ruudustiku uusi aspekte, mis võivad erineda teie ruudustikust. Alustage "KÜSIMUS:" ja ärge lisage midagi muud.} \\
            }
        }
    }
    & \\ \\

    \theutterance \stepcounter{utterance}  
    & & & \multicolumn{4}{p{0.6\linewidth}}{
        \cellcolor[rgb]{0.9,0.9,0.9}{
            \makecell[{{p{\linewidth}}}]{
                \texttt{\tiny{[GM$\langle$P2]}}
                \texttt{KÜSIMUS: Kas teie ruudustikus on muid X{-}i sümboleid, näiteks mõne teise veeru või rada kohta, mis ei kuulu mainitudud vertikaalsele või horisontaalsele mustrile?} \\
            }
        }
    }
     \\ \\

    \theutterance \stepcounter{utterance}  
    & & & \multicolumn{2}{p{0.3\linewidth}}{
        \cellcolor[rgb]{0.9,0.9,0.9}{
            \makecell[{{p{\linewidth}}}]{
                \texttt{\tiny{[GM$|$GM]}}
                \texttt{continue} \\
            }
        }
    }
    & & \\ \\

    \theutterance \stepcounter{utterance}  
    & & \multicolumn{4}{p{0.6\linewidth}}{
        \cellcolor[rgb]{0.9,0.9,0.9}{
            \makecell[{{p{\linewidth}}}]{
                \texttt{\tiny{[P1$\langle$GM]}}
                \texttt{VASTUS: Minu vertikaalne rida algab teisest reast (ruut 8). Seega meie ruudustikud ei ole täpselt samad, kuna minu vertikaalne rida on teie vertikaalse rida, kuid minu esimene rida (ruut 3) on tühi, samas kui teie kirjeldus viitab sellele, et teie vertikaalne rida algab teisest reast (mis ongi sama). Oota, kontrollin uuesti: minu kirjeldus ütles "ruudud 3, 8, 13, 18", mis tähendab esimest rida, teist, kolmandat ja neljandat. Teie kirjeldus aga ütles, et see algab teisest reast. Seega minu ruudustik on teie ruudustikust erinev, sest minu vertikaalne rida on üks ruut pikem (algab ruudust 3).} \\
\texttt{KÜSIMUS: Kas teie ruudustikus on muid X{-}i sümboleid, näiteks mõne teise veeru või rada kohta, mis ei kuulu mainitudud vertikaalsele või horisontaalsele mustrile?Alustage oma vastust "VASTUS:" ja ärge lisage midagi muud.} \\
            }
        }
    }
    & \\ \\

    \theutterance \stepcounter{utterance}  
    & \multicolumn{4}{p{0.6\linewidth}}{
        \cellcolor[rgb]{0.9,0.9,0.9}{
            \makecell[{{p{\linewidth}}}]{
                \texttt{\tiny{[P1$\rangle$GM]}}
                \texttt{VASTUS: Ei, minu ruudustikus ei ole muid X{-}i sümboleid. Kõik minu X{-}id kuuluvad vaid mainitud vertikaalsele (r2c3, r3c3, r4c3, r5c3) ja horisontaalsele (r5c1, r5c2, r5c3, r5c4, r5c5) mustrile.} \\
            }
        }
    }
    & & \\ \\

    \theutterance \stepcounter{utterance}  
    & & & \multicolumn{2}{p{0.3\linewidth}}{
        \cellcolor[rgb]{0.9,0.9,0.9}{
            \makecell[{{p{\linewidth}}}]{
                \texttt{\tiny{[GM$|$GM]}}
                \texttt{continue} \\
            }
        }
    }
    & & \\ \\

    \theutterance \stepcounter{utterance}  
    & & \multicolumn{4}{p{0.6\linewidth}}{
        \cellcolor[rgb]{0.9,0.9,0.9}{
            \makecell[{{p{\linewidth}}}]{
                \texttt{\tiny{[P1$\langle$GM]}}
                \texttt{Nüüd esitage küsimus, et teada saada minu ruudustiku uusi aspekte, mis võivad erineda teie ruudustikust. Alustage "KÜSIMUS:" ja ärge lisage midagi muud.} \\
            }
        }
    }
    & \\ \\

    \theutterance \stepcounter{utterance}  
    & \multicolumn{4}{p{0.6\linewidth}}{
        \cellcolor[rgb]{0.9,0.9,0.9}{
            \makecell[{{p{\linewidth}}}]{
                \texttt{\tiny{[P1$\rangle$GM]}}
                \texttt{KÜSIMUS: Kas teie horisontaalsel ridal (viimane rida) on kõik viis ruutu täidetud X{-}iga, või on seal mõni tühi ruut?} \\
            }
        }
    }
    & & \\ \\

    \theutterance \stepcounter{utterance}  
    & & & \multicolumn{2}{p{0.3\linewidth}}{
        \cellcolor[rgb]{0.9,0.9,0.9}{
            \makecell[{{p{\linewidth}}}]{
                \texttt{\tiny{[GM$|$GM]}}
                \texttt{continue} \\
            }
        }
    }
    & & \\ \\

    \theutterance \stepcounter{utterance}  
    & & \multicolumn{4}{p{0.6\linewidth}}{
        \cellcolor[rgb]{0.9,0.9,0.9}{
            \makecell[{{p{\linewidth}}}]{
                \texttt{\tiny{[GM$\rangle$P2]}}
                \texttt{VASTUS: Ei, minu ruudustikus ei ole muid X{-}i sümboleid. Kõik minu X{-}id kuuluvad vaid mainitud vertikaalsele (r2c3, r3c3, r4c3, r5c3) ja horisontaalsele (r5c1, r5c2, r5c3, r5c4, r5c5) mustrile.} \\
\texttt{KÜSIMUS: Kas teie horisontaalsel ridal (viimane rida) on kõik viis ruutu täidetud X{-}iga, või on seal mõni tühi ruut?Alustage oma vastust "VASTUS:" ja ärge lisage midagi muud.} \\
            }
        }
    }
    & \\ \\

    \theutterance \stepcounter{utterance}  
    & & & \multicolumn{4}{p{0.6\linewidth}}{
        \cellcolor[rgb]{0.9,0.9,0.9}{
            \makecell[{{p{\linewidth}}}]{
                \texttt{\tiny{[GM$\langle$P2]}}
                \texttt{VASTUS: Minu horisontaalsel ridal (viimane rida) on kõik viis ruutu täidetud X{-}iga.} \\
            }
        }
    }
     \\ \\

    \theutterance \stepcounter{utterance}  
    & & & \multicolumn{2}{p{0.3\linewidth}}{
        \cellcolor[rgb]{0.9,0.9,0.9}{
            \makecell[{{p{\linewidth}}}]{
                \texttt{\tiny{[GM$|$GM]}}
                \texttt{continue} \\
            }
        }
    }
    & & \\ \\

    \theutterance \stepcounter{utterance}  
    & & \multicolumn{4}{p{0.6\linewidth}}{
        \cellcolor[rgb]{0.9,0.9,0.9}{
            \makecell[{{p{\linewidth}}}]{
                \texttt{\tiny{[GM$\rangle$P2]}}
                \texttt{Nüüd esitage küsimus, et teada saada minu ruudustiku uusi aspekte, mis võivad erineda teie ruudustikust. Alustage "KÜSIMUS:" ja ärge lisage midagi muud.} \\
            }
        }
    }
    & \\ \\

    \theutterance \stepcounter{utterance}  
    & & & \multicolumn{4}{p{0.6\linewidth}}{
        \cellcolor[rgb]{0.9,0.9,0.9}{
            \makecell[{{p{\linewidth}}}]{
                \texttt{\tiny{[GM$\langle$P2]}}
                \texttt{KÜSIMUS: Kas teie ruudustiku esimene rida (r1) on täiesti tühi, või sisaldab see muid X{-}i sümboleid?} \\
            }
        }
    }
     \\ \\

    \theutterance \stepcounter{utterance}  
    & & & \multicolumn{2}{p{0.3\linewidth}}{
        \cellcolor[rgb]{0.9,0.9,0.9}{
            \makecell[{{p{\linewidth}}}]{
                \texttt{\tiny{[GM$|$GM]}}
                \texttt{continue} \\
            }
        }
    }
    & & \\ \\

    \theutterance \stepcounter{utterance}  
    & & \multicolumn{4}{p{0.6\linewidth}}{
        \cellcolor[rgb]{0.9,0.9,0.9}{
            \makecell[{{p{\linewidth}}}]{
                \texttt{\tiny{[P1$\langle$GM]}}
                \texttt{VASTUS: Minu horisontaalsel ridal (viimane rida) on kõik viis ruutu täidetud X{-}iga.} \\
\texttt{KÜSIMUS: Kas teie ruudustiku esimene rida (r1) on täiesti tühi, või sisaldab see muid X{-}i sümboleid?Alustage oma vastust "VASTUS:" ja ärge lisage midagi muud.} \\
            }
        }
    }
    & \\ \\

    \theutterance \stepcounter{utterance}  
    & \multicolumn{4}{p{0.6\linewidth}}{
        \cellcolor[rgb]{0.9,0.9,0.9}{
            \makecell[{{p{\linewidth}}}]{
                \texttt{\tiny{[P1$\rangle$GM]}}
                \texttt{VASTUS: Minu ruudustiku esimene rida (r1) on täiesti tühi.} \\
            }
        }
    }
    & & \\ \\

    \theutterance \stepcounter{utterance}  
    & & & \multicolumn{2}{p{0.3\linewidth}}{
        \cellcolor[rgb]{0.9,0.9,0.9}{
            \makecell[{{p{\linewidth}}}]{
                \texttt{\tiny{[GM$|$GM]}}
                \texttt{continue} \\
            }
        }
    }
    & & \\ \\

    \theutterance \stepcounter{utterance}  
    & & \multicolumn{4}{p{0.6\linewidth}}{
        \cellcolor[rgb]{0.9,0.9,0.9}{
            \makecell[{{p{\linewidth}}}]{
                \texttt{\tiny{[P1$\langle$GM]}}
                \texttt{Nüüd esitage küsimus, et teada saada minu ruudustiku uusi aspekte, mis võivad erineda teie ruudustikust. Alustage "KÜSIMUS:" ja ärge lisage midagi muud.} \\
            }
        }
    }
    & \\ \\

    \theutterance \stepcounter{utterance}  
    & \multicolumn{4}{p{0.6\linewidth}}{
        \cellcolor[rgb]{0.9,0.9,0.9}{
            \makecell[{{p{\linewidth}}}]{
                \texttt{\tiny{[P1$\rangle$GM]}}
                \texttt{KÜSIMUS: Kas teie ruudustikus on muid X{-}i sümboleid, mis ei kuulu mainitudud vertikaalsele või horisontaalsele mustrile, näiteks mõnes teises veerus või reas?} \\
            }
        }
    }
    & & \\ \\

    \theutterance \stepcounter{utterance}  
    & & & \multicolumn{2}{p{0.3\linewidth}}{
        \cellcolor[rgb]{0.9,0.9,0.9}{
            \makecell[{{p{\linewidth}}}]{
                \texttt{\tiny{[GM$|$GM]}}
                \texttt{continue} \\
            }
        }
    }
    & & \\ \\

    \theutterance \stepcounter{utterance}  
    & & \multicolumn{4}{p{0.6\linewidth}}{
        \cellcolor[rgb]{0.9,0.9,0.9}{
            \makecell[{{p{\linewidth}}}]{
                \texttt{\tiny{[GM$\rangle$P2]}}
                \texttt{VASTUS: Minu ruudustiku esimene rida (r1) on täiesti tühi.} \\
\texttt{KÜSIMUS: Kas teie ruudustikus on muid X{-}i sümboleid, mis ei kuulu mainitudud vertikaalsele või horisontaalsele mustrile, näiteks mõnes teises veerus või reas?Alustage oma vastust "VASTUS:" ja ärge lisage midagi muud.} \\
            }
        }
    }
    & \\ \\

    \theutterance \stepcounter{utterance}  
    & & & \multicolumn{4}{p{0.6\linewidth}}{
        \cellcolor[rgb]{0.9,0.9,0.9}{
            \makecell[{{p{\linewidth}}}]{
                \texttt{\tiny{[GM$\langle$P2]}}
                \texttt{VASTUS: Ei, minu ruudustikus ei ole muid X{-}i sümboleid peale minu nimetatud vertikaalse ja horisontaalse rida.} \\
            }
        }
    }
     \\ \\

    \theutterance \stepcounter{utterance}  
    & & & \multicolumn{2}{p{0.3\linewidth}}{
        \cellcolor[rgb]{0.9,0.9,0.9}{
            \makecell[{{p{\linewidth}}}]{
                \texttt{\tiny{[GM$|$GM]}}
                \texttt{continue} \\
            }
        }
    }
    & & \\ \\

    \theutterance \stepcounter{utterance}  
    & & \multicolumn{4}{p{0.6\linewidth}}{
        \cellcolor[rgb]{0.9,0.9,0.9}{
            \makecell[{{p{\linewidth}}}]{
                \texttt{\tiny{[GM$\rangle$P2]}}
                \texttt{Nüüd tee otsus. Mida sa arvad: kas sinu ruudustik ja teine kirjeldatud ruudustik on samad?  Kirjuta "OTSUS: sama ruudustik", kui arvad, et need on sama ruudustik, või "OTSUS: erinevad ruudustikud", kui arvad, et need on erinevad ruudustikud. Ära lisa midagi muud.} \\
            }
        }
    }
    & \\ \\

    \theutterance \stepcounter{utterance}  
    & & & \multicolumn{4}{p{0.6\linewidth}}{
        \cellcolor[rgb]{0.9,0.9,0.9}{
            \makecell[{{p{\linewidth}}}]{
                \texttt{\tiny{[GM$\langle$P2]}}
                \texttt{OTSUS: erinevad ruudustikud} \\
            }
        }
    }
     \\ \\

    \theutterance \stepcounter{utterance}  
    & & & \multicolumn{2}{p{0.3\linewidth}}{
        \cellcolor[rgb]{0.9,0.9,0.9}{
            \makecell[{{p{\linewidth}}}]{
                \texttt{\tiny{[GM$|$GM]}}
                \texttt{continue} \\
            }
        }
    }
    & & \\ \\

    \theutterance \stepcounter{utterance}  
    & & & \multicolumn{2}{p{0.3\linewidth}}{
        \cellcolor[rgb]{0.9,0.9,0.9}{
            \makecell[{{p{\linewidth}}}]{
                \texttt{\tiny{[GM$|$GM]}}
                \texttt{loss} \\
            }
        }
    }
    & & \\ \\

    \theutterance \stepcounter{utterance}  
    & & \multicolumn{4}{p{0.6\linewidth}}{
        \cellcolor[rgb]{0.9,0.9,0.9}{
            \makecell[{{p{\linewidth}}}]{
                \texttt{\tiny{[P1$\langle$GM]}}
                \texttt{VASTUS: Ei, minu ruudustikus ei ole muid X{-}i sümboleid peale minu nimetatud vertikaalse ja horisontaalse rida.} \\
\texttt{Nüüd tee otsus. Mida sa arvad: kas sinu ruudustik ja teine kirjeldatud ruudustik on samad?  Kirjuta "OTSUS: sama ruudustik", kui arvad, et need on sama ruudustik, või "OTSUS: erinevad ruudustikud", kui arvad, et need on erinevad ruudustikud. Ära lisa midagi muud.} \\
            }
        }
    }
    & \\ \\

    \theutterance \stepcounter{utterance}  
    & \multicolumn{4}{p{0.6\linewidth}}{
        \cellcolor[rgb]{0.9,0.9,0.9}{
            \makecell[{{p{\linewidth}}}]{
                \texttt{\tiny{[P1$\rangle$GM]}}
                \texttt{OTSUS: erinevad ruudustikud} \\
            }
        }
    }
    & & \\ \\

    \theutterance \stepcounter{utterance}  
    & & & \multicolumn{2}{p{0.3\linewidth}}{
        \cellcolor[rgb]{0.9,0.9,0.9}{
            \makecell[{{p{\linewidth}}}]{
                \texttt{\tiny{[GM$|$GM]}}
                \texttt{continue} \\
            }
        }
    }
    & & \\ \\

    \theutterance \stepcounter{utterance}  
    & & & \multicolumn{2}{p{0.3\linewidth}}{
        \cellcolor[rgb]{0.9,0.9,0.9}{
            \makecell[{{p{\linewidth}}}]{
                \texttt{\tiny{[GM$|$GM]}}
                \texttt{loss} \\
            }
        }
    }
    & & \\ \\

\end{supertabular}
}

\end{document}
